\section{Literature Review}

A lot of research has already been done on the differences between client-side rendering and server-side rendering within web applications. In most studies, the focus is on performance in terms of speed and user experience. Nordström and Dixelius (2023) show, for example, that SSR can deliver better load performance than CSR for larger and more complex web pages, while for smaller pages the difference is minimal. Clapton et al. (2025) also show that SSR scores better on metrics such as Time To Interactive (TTI) and Largest Contentful Paint (LCP), especially with larger data volumes.

Ojala (2025) describes that SSR offers advantages in terms of search engine optimization (SEO) and initial load time. At the same time, SSR introduces extra complexity in the development process, for example through server load and the use of hydration mechanisms. What stands out to me in these studies is that performance is almost always interpreted as speed or user experience, and less often as efficient use of resources.

In some studies, CPU and memory usage are also examined. Gaddam (2025) reports that SSR on average leads to lower CPU and memory load during the render phase, particularly on client devices. This is explained by the fact that a larger part of the processing takes place on the server, meaning the client needs to execute less JavaScript. Clapton et al. (2025) also analyze the use of the JavaScript Heap and show that with larger datasets, the client-side load with CSR increases significantly compared to SSR. This suggests that the choice of render strategy not only affects speed, but also the load on the user's device.

When looking at the energy consumption aspect of web applications, it appears that this topic is less directly linked to specific render strategies. Thiagarajan et al. (2012) show with a hardware-based measurement setup that JavaScript and CSS processing in the browser have a significant influence on the energy consumption of mobile devices. Recent studies, such as Kalliola and Vepsäläinen (2025), indicate that measuring energy consumption of web applications is complex. They state that CPU usage can serve as a useful indicator for energy consumption, but that existing online ``energy calculators'' often give inaccurate or inconsistent results.

What stands out in the existing literature is that resource usage and energy efficiency rarely form the primary research objective. CPU and memory measurements are often included as supporting metrics in performance analyses, but not as central evaluation criteria.

Based on this literature, it can be concluded that SSR and CSR have been extensively compared in terms of speed, user experience, and SEO. It has also been shown that render strategies influence CPU and memory usage on client devices. However, there appears to be limited research that explicitly compares SSR and CSR with resource efficiency as the focus, and specifically within a mobile context. This forms the motivation for the current research, in which the emphasis shifts from ``performance as speed'' to ``performance as resource efficiency''.
